% % 01-simu-enunciado2.tex
% % Simulación de sucesos discretos y optimización

% Consideremos un almacén de dos productos cuyos precios de venta al público son de 2.5 y 3.5 euros la unidad, respectivamente. La llegada de clientes al almacén se distribuye según un proceso de Poisson de parámetro $\lambda = 1.5$ clientes por hora, y la cantidad de productos demandados por cada uno de ellos tiene la siguiente distribución:
% \begin{table}[ht!]
%     \centering
%     \begin{tabular}{|c|c|c|c|c|}\hline
%         \rowcolor{colorA}\emphWhite{Demanda} & \emphWhite{1 unidad} & \emphWhite{2 unidades} & \emphWhite{3 unidades} & \emphWhite{4 unidades} \\ \hline
%         Producto 1 (\euro 2.5)               & 0.3                  & 0.4                    & 0.2                    & 0.1                    \\ \hline
%         Producto 2 (\euro 3.5)               & 0.2                  & 0.2                    & 0.4                    & 0.2                    \\ \hline
%     \end{tabular}
% \end{table}

% Para satisfacer la demanda de sus clientes, el dueño del almacén mantiene un stock de productos. La política de pedidos al distribuidor es periódica, es decir, todos los viernes a primera hora realiza un pedido en el que, tomando como referencia el nivel de inventario de los productos en ese momento, se solicitan las unidades necesarias para que el nivel del inventario de cada producto llegue a 1000 y 1500 unidades, respectivamente en cada producto.

% Asociado a cada pedido realizado al proveedor, existe un coste fijo (coste de preparación) de 100 euros, independientemente de las unidades demandadas. Adicionalmente, el coste por unidad incluida en el pedido depende de la cantidad solicitada, habiendo descuentos por cantidad. Si el número de unidades demandadas del primer producto es menor o igual a 600, el precio es de 1 euro la unidad; mientras que si se piden más de 600, el precio desciende a 0.75 euros. Para el segundo producto, si el número de unidades demandadas es menor que 800, el precio es de 1.5 euros la unidad; mientras que si se piden más de 800, el precio desciende a 1.25 euros.

% El tiempo que tarda en ser servido el pedido por los proveedores (tiempo líder) sigue una distribución normal de media 48 horas y desviación típica 3.5 horas, pagándose en ese momento. Se ha llegado a un acuerdo con el proveedor de forma que, tomando como referencia las 48 horas que tarda en media un pedido en ser servido, si el pedido llega con 3 horas de retraso en la entrega, se realiza un descuento del 0.03\% del valor del pedido, encareciéndose en la misma cantidad en el caso de que el pedido llegue con al menos 3 horas de adelanto.

% El dueño del almacén debe pagar 0.0002 euros por unidad del producto y unidad de tiempo, asociado al almacenamiento físico de los productos (alquiler del local, refrigeración, etc.). En el caso de que al llegar un cliente éste solicite una cantidad mayor que la que hay en inventario, se le sirve lo que queda, perdiendo la venta restante.

% \begin{enumerate}
%     \item Simular el comportamiento del almacén durante un periodo de tiempo de 5 meses para estimar:
%           \begin{itemize}
%               \item El beneficio esperado,
%               \item La proporción de clientes cuya demanda se satisface completamente,
%               \item El porcentaje de tiempo que el nivel del inventario permanece a cero.
%           \end{itemize}
%           Para ello, supondremos que el nivel del inventario inicial es de 70 unidades de ambos productos.

%     \item Representar gráficamente la evolución del nivel del inventario durante los 5 meses y durante los 5 primeros días.

%     \item Mediante el uso de la metaheurística de \textit{recocido simulado}, identificar cuál será la política de pedidos óptima. Es decir:
%           \begin{itemize}
%               \item Identificar cada cuánto tiempo se deberían realizar los pedidos,
%               \item El valor de referencia del inventario para identificar el número de unidades a solicitar en la política de pedidos periódica.
%           \end{itemize}
% \end{enumerate}

% \noindent \textbf{Nota:} El almacén permanece abierto las 24 horas al día de lunes a domingo. El proveedor y la empresa transportista (que sirve los pedidos) también trabajan las 24 horas, sin cerrar en ningún momento.

\begin{figure}[ht!]
    \centering
    \begin{tikzpicture}
        % \draw[fill=GithubGreen, draw=none, opacity=0.5] (-0.5,-.5) rectangle (7.5,2);
        % \draw[fill=black, text=white,draw=none] (7.75,-.5) rectangle (15.5,2);
        % \draw[fill=colorE, draw=none, opacity=0.95] (-0.5,-2.5) rectangle (7.5,-.75);
        % \node[below right] (State) at (0,1.8){\grotesk\large Variables de Estado};
        % \node[below right] (State2) at (0,1.25) {$x$: Cantidad de productos en inventario};
        % \node[below right] (State3) at (0,0.65) {$y$: Cantidad de productos pedidos al proveedor};
        % \node[below right,text=white] (Counters) at (7.8,1.9){\grotesk\large Contadores};
        % \node[below right,text=white] (Counters2) at (7.8,1.3) {$C$: Costo total de pedidos en tiempo $t$};
        % \node[below right,text=white] (Counters3) at (7.8,0.7) {$H$: Costo total de almacenamiento en tiempo $t$};
        % \node[below right] (Counters3) at (7.8,0.1) {$R$: Beneficio total al finalizar $t$};
        % \node[below right] (Event) at (0,-0.75){\grotesk\large Eventos};
        % \node[below right] (Event1) at (0,-1.25) {$t_0$: Tiempo de llegada de siguiente cliente};
        % \node[below right] (Event2) at (0,-1.85) {$t_1$: Tiempo de llegada de siguiente pedido};
        % Box
        \draw[fill=black, draw=none, opacity=0.5] (4,0) rectangle (12,10);
        \node[text=white,right] (Holding) at (4,9.5){\Large\grotesk Almacén};
        \node[text=white,below right=1mm of Holding, align=left] (Holding1){\large\ttfamily Costo de almacenamiento \\($h = 0.0002 \times u \times t$)};
        % Product red
        \draw[-, thick,colorE] (3,4)node[left]{\Large $P_1$ (\euro 2.5)} -- ++(1,0);
        \node[colorE,below] (p1Max) at (3,3.8){\ttfamily Max=1000};
        \draw[fill=colorE, draw=none] (4,0) rectangle (8,4); % Bottom
        \draw [very thick, draw=white] (4,0) grid  (8,4); % Bottom
        % Product Blue
        \draw[-, thick,colorA] (13,8)node[right]{\Large $P_2$ (\euro 3.5)} -- ++(-1,0);
        \node[colorA,below] (p1Max) at (13,7.8){\ttfamily Max=1500};
        \draw[fill=colorA, draw=none] (8,0) rectangle (12,8); % Bottom
        \draw [very thick, draw=white] (8,0) grid  (12,8); % Bottom
        % 
        % Indicadores p y P
        % \draw[-, thick,gray] (3,2)node[left]{\Large $p$} -- ++(1,0);
        % Llegadas
        \draw[-Triangle, line width=0.7mm, GithubGreen] (\textwidth,11) -- ++(-5.5,0)node[above,pos=0.6]{\bfseries\small PROCESO DE LLEGADAS} -- ++(0,-1.5);
        \node[blackLabel, left] (productB) at (16,11){\grotesk Clientes};
        \node[greenLabel,above left=6mm of productB, text width=4.6cm] (a) {\ttfamily Demanda de productos};
        \node[greenLabel,above =1mm of a, text width=4.6cm] (b) {\ttfamily Distribución discreta};
        \node[greenLabel,above =1mm of b, text width=4.6cm] (c) {\ttfamily Precio de venta (r)};
        % Pedidos
        \draw[-Triangle, line width=0.7mm, colorD] (2,11) -- ++(5.5,0)node[above,pos=0.6]{\bfseries\small REALIZACIÓN DE PEDIDO} -- ++(0,-1.5);
        \node[blackLabel, text=white,right] (productA) at (0,11){\grotesk Proveedor};
        \node[bLabel,above right=6mm of productA, text width=4.6cm] (a) {\ttfamily Costo base ($K=$ 100 \euro)};
        \node[bLabel,above =1mm of a, text width=4.6cm] (b) {\ttfamily Coste por unidad};
        \node[bLabel,above =1mm of b, text width=4.6cm] (c) {\ttfamily Descuento por cantidad};

        % Pedidos a Almacen:
        \draw[-Triangle, line width=0.6mm, dashed, colorD] (1.2,10.5) -- ++(0,-6) -- ++(2,0);
        \node[bLabel,below,align=left] (politic) at (1.2,8){\bfseries\small Llegada de pedido\\ Tiempo líder ($L$)\\ $\mathcal{N}(48,3.5) $};
        \node[text=gray, align=left] (Holding2) at (1,2){\ttfamily Política de pedidos:\\ periódica\\ Viernes a primera hora};
        \draw[line width=0.5mm,-,black] (0,0) -- (\textwidth,0);
        % \draw[line width=0.5mm,-,red] (0,0) -- (0,\textheight);
        \node[shape=rectangle,fill=GithubGreen!20, below right] (somenode) at (0,-0.2) {
            \begin{tabular}{ccccc}\toprule
                Demanda    & 1 unidad & 2 unidades & 3 unidades & 4 unidades \\ \hline
                Producto 1 & 0.3      & 0.4        & 0.2        & 0.1        \\ \hline
                Producto 2 & 0.2      & 0.2        & 0.4        & 0.2        \\ \hline
                \bottomrule
            \end{tabular}
        };
    \end{tikzpicture}
    % \caption{asdkj}
\end{figure}
\nocite{*}
% \begin{enumerate}
% \item Gráfico de $p(i)$ vs iteraciones. Para ver el valor de la probabilidad de aceptación de soluciones peores (0.9)
% \item Gráfico de soluciones vs iteraciones.
% \item Gráfico de $f^*$ vs iteraciones.
% \item Registrar el costo de la venta perdida*
% \end{enumerate}
% https://ucreanop.com/wp-content/uploads/2021/01/Capitulo-16.pdf